\resume
\selectlanguage{french}
\vspace*{-6cm}
\begin{abstract}
Ce mémoire porte sur la mise en place d'un système de détection et de prévention d'intrusion (IDS/IPS) pour sécuriser l'infrastructure réseau du Centre National Hospitalier et Universitaire Hubert Koutoukou Maga (CNHU-HKM) de Cotonou. Un audit réalisé lors d'un stage au service informatique et réseau a montré que le segment analysé ne présente ni segmentation visible ni mécanisme de détection d'intrusion. Après une étude comparative des solutions, le moteur Suricata est retenu pour un déploiement sur la plateforme de pare-feu OPNsense, prévu d'abord en mode détection puis en mode prévention, avec des règles personnalisées pour les flux DICOM et HL7. Un protocole d'évaluation dans un environnement virtualisé distinct de la production, fondé sur sept scénarios d'attaque rattachés au référentiel MITRE ATT\&CK et rejoués dans trois états (sans détection, IDS, IPS), est entièrement défini avec sa grille de mesures et les preuves à recueillir. L'expérimentation n'ayant pas été exécutée dans le cadre de ce mémoire, aucun résultat mesuré n'est présenté et aucune performance n'est affirmée sans mesure.

\paragraph{}
\textbf{Mots clés}: IDS/IPS, détection d'intrusion, Suricata, OPNsense, CNHU-HKM, sécurité réseau.
\end{abstract}

\newpage
\thispagestyle{empty}
\selectlanguage{english}
\addcontentsline{toc}{chapter}{Abstract}
\begin{abstract}

This thesis deals with the deployment of an intrusion detection and prevention system (IDS/IPS) to secure the network infrastructure of the Hubert Koutoukou Maga National Teaching Hospital (CNHU-HKM) in Cotonou. An audit carried out during an internship in the hospital's IT and network department showed that the analysed segment has neither visible segmentation nor any intrusion detection mechanism. After a comparative study of the available solutions, the Suricata engine is selected for deployment on the OPNsense firewall platform, planned first in detection mode and then in prevention mode, with custom rules for DICOM and HL7 flows. An evaluation protocol in a virtualised environment separate from production, based on seven attack scenarios mapped to the MITRE ATT\&CK framework and replayed in three states (no detection, IDS, IPS), is fully defined together with its measurement grid and the evidence to be collected. As the experiment was not carried out within this work, no measured results are reported and no performance claim is made without measurement.

\paragraph{}
\textbf{Key words}: IDS/IPS, intrusion detection, Suricata, OPNsense, CNHU-HKM, network security.
\end{abstract}
